% BEGIN PROFESSOR FEEDBACK PF17: Implementation - Result artifacts
\subsection{Result artifacts}

Each run has its own directory that is never overwritten, so results from
different runs cannot be mixed, and a run may only create a file that does not yet
exist. Prepared edges, fingerprints, and test-set predictions are stored as
Parquet, the protocol and the test summaries as JavaScript Object Notation
(JSON), the chapter tables as comma-separated values (CSV),
and the figures as PDF. Output files are first fully written and synchronised in
a private attempt directory. Files are committed without replacement, and a
multi-file stage is published by an atomic directory rename after its manifest is
written last. The public \texttt{features}, \texttt{metrics}, and related
directories contain hard links to verified stage payloads, preserving the
filenames consumed by the Results chapter.

Preparation and execution have separate resume rules. During Real-CATS
preparation, the SQLite index records each completed address and can be resumed
only before the experiment settings are fixed. During execution, every attempt
has its own identifier and scratch directory. HF commits one subgraph-count
checkpoint for every ordering and window. Real-CATS commits one for every layer, ordering, and
window. A new attempt may reuse earlier valid commits. A completed run rejects
both a fresh execution and a resume.

The integrity chain is acyclic. The locked protocol and prepared data identify
the execution definition. These identify every checkpoint and stage. The
result inventory hashes the committed stage manifests and public files, and the
execution manifest records the inventory hash. Each inventory and manifest has
a separate checksum file and does not include itself. The final
\texttt{verify} command checks this chain and the scientific schemas without
writing to the run. Ordinary verification distinguishes a valid non-citable
development or smoke run from corruption. Final thesis acceptance additionally
requires \texttt{--require-citable}.

The only appended file is \texttt{run\_log.jsonl}, which stores one JSON record
per line in JSON Lines (JSONL) format. Its writes are thread-safe
and synchronised to disk. A native operation that runs for several minutes
emits a heartbeat every five minutes with its elapsed time, allocated threads,
free disk, and latest committed checkpoint. The output layout is shown below.

\begin{lstlisting}[caption={Locked Results artifact layout.},
  label={lst:filenames}]
results/experiments/<run-id>/
  preparation_state.json
  protocol.json / protocol.sha256
  execution_definition.json / execution_definition.sha256
  attempts/<attempt-id>/
    started.json, completed.json or failed.json
    scratch/duckdb/
    staging/
  checkpoints/<checkpoint-id>/
    counts.parquet
    metadata.json
  stages/<stage-id>/
    payload/
    stage_manifest.json
  run_log.jsonl
  data/
  features/
  predictions/
  metrics/
  figures/
  tables/
  result_artifacts.json / result_artifacts.sha256
  execution_manifest.json / execution_manifest.sha256
\end{lstlisting}
% END PROFESSOR FEEDBACK PF17

